\chapter{Background and Related Work}\label{sec:literature-review}

\paragraph{Scope relative to the canonical model.} Unless stated otherwise, the relay-processing methods reviewed in this chapter are discussed in the context of the canonical memoryless channel model adopted throughout Chapter~\ref{sec:experiments}. Channel-memory effects, sequence detection, and channel models with intersymbol interference (ISI) are intentionally deferred to the unknown-channel study of Chapter~\ref{sec:unknown-channels}.

\section{Classical Relay Strategies}\label{sec:classical-relay-strategies}

The fundamental question that motivates this thesis is: can a learned relay function \(f_\theta(\cdot)\) outperform the classical relay functions (AF scaling, DF regeneration) by exploiting patterns in the received signal that analytical methods do not capture? This question is particularly relevant at low SNR, where both AF (noise amplification) and DF (decoding errors) have well-known limitations, and where the non-linear decision boundary learned by a neural network may provide a potential advantage.

\subsection{Amplify-and-Forward (AF)}\label{sec:classical-relay-strategies-af}
The AF relay amplifies the received signal by a gain factor \(G\) that normalizes the output power: (Equation~\ref{eq:af-gain})
\addcontentsline{equ}{listequations}{\protect\numberline{\theequation}\quad af-gain}

\begin{equation}\protect\phantomsection\label{eq:af-gain}{G = \sqrt{\frac{P_{\text{target}}}{\mathbb{E}[|y_R|^2]}}
}\end{equation}

The end-to-end SNR for AF relay is the standard harmonic-mean expression \cite{Laneman2004CooperativeDiversity,TseViswanath2005Fundamentals}: (Equation~\ref{eq:af-snr})

\begin{equation}\protect\phantomsection\label{eq:af-snr}{\text{SNR}_{\text{eff}}^{\text{AF}} = \frac{\text{SNR}_1 \cdot \text{SNR}_2}{\text{SNR}_1 + \text{SNR}_2 + 1}
}\end{equation}
\addcontentsline{equ}{listequations}{\protect\numberline{\theequation}\quad af-snr}

This form assumes \emph{CSI-assisted variable-gain} AF (Laneman et al.), in which the relay computes a fresh gain from its own \emph{instantaneous} first-hop channel state for every symbol. That is a different scheme from Equation~\ref{eq:af-gain} as implemented in this thesis: the simulator's gain $G$ is a single scalar, estimated once per simulated block from the empirical average $\mathbb{E}[|y_R|^2]$ over that block rather than from a per-symbol channel estimate, which is the power-normalized \emph{fixed-gain} variant. Equations~\ref{eq:af-snr}--\ref{eq:af-p-eff} are therefore standard background from the cooperative-diversity literature, included to state the well-known noise-amplification limitation of AF; they are not the closed form for, and are not compared numerically against, the Monte Carlo AF results reported elsewhere in this thesis, which come entirely from simulation. The additive $+1$ in the denominator of Equation~\ref{eq:af-snr} originates from the relay's own noise being re-amplified along with the signal under instantaneous variable-gain scaling; the fixed-gain variant implemented here has a different closed form, not derived here since it is not needed for anything measured.

This expression reveals a fundamental limitation: even when one hop has high SNR, the effective SNR is bottlenecked by the weaker hop. Moreover, the relay amplifies both signal and noise from the first hop, leading to noise accumulation at the destination.

\subsection{Decode-and-Forward (DF)}\label{sec:classical-relay-strategies-df}
The DF relay demodulates the received signal, recovers the transmitted bits, and re-modulates clean symbols:

\begin{equation}\protect\phantomsection\label{eq:df-demod-remod}{\hat{b}_R = \text{demod}(y_R), \quad x_R = \text{mod}(\hat{b}_R)
}\end{equation}
\addcontentsline{equ}{listequations}{\protect\numberline{\theequation}\quad df-demod-remod}

The end-to-end BER for DF is: (Equation~\ref{eq:df-ber-hops})

\begin{equation}\protect\phantomsection\label{eq:df-ber-hops}{P_e^{\text{DF}} = P_{e,1}(1 - P_{e,2}) + (1 - P_{e,1}) \cdot P_{e,2}
}\end{equation}
\addcontentsline{equ}{listequations}{\protect\numberline{\theequation}\quad df-ber-hops}

where \(P_{e,1}\) and \(P_{e,2}\) are the BER of the first and second hops, respectively: an end-to-end bit error occurs exactly when \emph{one} of the two hops flips the bit (two flips cancel), so this expands to \(P_{e,1} + P_{e,2} - 2P_{e,1}P_{e,2}\), the same expression used in the AWGN analysis (Equation~\ref{eq:df-ber-awgn}). For BPSK modulation over AWGN (real-valued noise with variance \(1/\text{SNR}\)), each hop's BER is \(P_e = Q\left(\sqrt{\text{SNR}}\right)\). DF provides clean regeneration at high SNR but suffers from error propagation when the first-hop BER is non-negligible. As discussed in Remark~\ref{rem:df-terminology}, the ``DF'' baseline used here is the uncoded symbol-wise slicer, not practical finite-length coded block-DF; its optimality claims are therefore confined to the uncoded setting.

\subsection{Other Classical Relay Techniques}

AF and DF represent two canonical extremes of the relay processing spectrum: AF performs minimal processing (linear scaling) while DF performs maximal classical processing (full regeneration). Additional techniques (e.g. CF, EF, CoF) occupy intermediate points along this spectrum. In this thesis, the AI-based relay strategies are designed to \emph{learn} the relay processing function from data, potentially discovering strategies that subsume or outperform these fixed classical schemes. The focus on AF and DF as classical baselines is therefore deliberate: on the canonical matched-channel benchmark they bracket the classical design space and serve as representative reference baselines against which the learned strategies are evaluated; they should not be read as strict lower and upper performance bounds, since under channel mismatch a learned relay can outperform DF (Chapter~\ref{sec:unknown-channels}). These additional relay strategies are reviewed for completeness; they are not simulated in this thesis because the studied system assumes a two-hop link without a direct source-destination path.

\section{Machine Learning in Wireless Communication}\label{sec:machine-learning-in-wireless-communication}

The application of machine learning to physical-layer wireless communication has gained significant momentum in recent years, driven by the ability of neural networks to learn complex non-linear mappings directly from data. Deep learning has been applied to channel estimation \cite{YeLiJuang2018OFDMDL}, signal detection \cite{SamuelDiskinWunder2019MIMODetect}, autoencoder-based end-to-end communication \cite{Dorner2018OverTheAirDL}, and resource allocation \cite{Sun2018LearnToOptimize}. These approaches learn complex mappings from data, potentially outperforming hand-crafted algorithms in scenarios where analytical solutions are intractable or suboptimal.

\subsection{Theoretical Basis: Universal Approximation and Denoising}\label{sec:theoretical-basis-universal-approximation-and-denoising}

The theoretical justification for applying neural networks to relay signal processing rests on two pillars. First, the \textbf{universal approximation theorem} \cite{Cybenko1989UAT,Hornik1991UAT,GoodfellowBengioCourville2016DeepLearning} guarantees that a feedforward network with a single hidden layer and a non-linear activation function can approximate any continuous function on a compact domain to arbitrary accuracy, given sufficient width. For relay denoising, the target function maps noisy observations to clean transmitted symbols:

\begin{equation}\protect\phantomsection\label{eq:optimal-denoiser}{f^*: \mathbb{R}^{2w+1} \to [-1, 1], \quad f^*(y_{i-w}, \dots, y_{i+w}) = \mathbb{E}[x_i \mid y_{i-w}, \dots, y_{i+w}]
}\end{equation}
\addcontentsline{equ}{listequations}{\protect\numberline{\theequation}\quad optimal-denoiser}

This conditional expectation \(f^*\) is the Bayes-optimal denoiser: it minimizes the mean squared error (MSE) over all possible estimators. For BPSK symbols corrupted by AWGN, \(f^*\) reduces to the posterior mean:

\begin{equation}\protect\phantomsection\label{eq:optimal-denoiser-awgn}{f^*(y) = \tanh\left(\frac{y}{\sigma^2}\right)
}\end{equation}
\addcontentsline{equ}{listequations}{\protect\numberline{\theequation}\quad optimal-denoiser-awgn}

which is a smooth sigmoid approaching the hard-decision signum function as \(\sigma^2 \to 0\). A single tanh unit with the right input scaling represents this posterior \emph{exactly} for a fixed \(\sigma^2\). The deployed relay, however, is trained at SNR \(\in\{5,10,15\}\)~dB and evaluated over \(0\)--\(20\)~dB, so one set of weights must approximate a whole \emph{family} of posteriors indexed by \(\sigma^2\), and on the fading channel a mixture over random effective SNRs. That a 169-parameter network approximates that family closely enough to track the classical baselines is an empirical finding of this thesis, not a consequence of the exact single-\(\sigma^2\) representation.

Second, the \textbf{bias-variance decomposition} provides a framework for understanding model complexity:

\begin{equation}\protect\phantomsection\label{eq:bias-variance-decomp}{\mathbb{E}[(\hat{x} - x)^2] = \text{Bias}^2(\hat{x}) + \text{Var}(\hat{x}) + \sigma^2_{\text{irreducible}}
}\end{equation}
\addcontentsline{equ}{listequations}{\protect\numberline{\theequation}\quad bias-variance-decomp}

The irreducible term is the minimum achievable MSE, set by the channel noise. Adding parameters reduces bias but increases variance; since the target \(f^*\) here is essentially a soft threshold, the bias term is already small for modest networks, so on this reasoning additional capacity would buy little. This is the standard argument for preferring a small model on a simple target. Whether it is what actually happens in the measurements of this thesis is a separate question, and Section~\ref{sec:capacity-and-overfitting} is careful to report only that larger evaluated models did not improve BER, without attributing that to overfitting.

\subsection{Prior Work in Deep Learning for Physical-Layer Processing}\label{sec:prior-work-in-deep-learning-for-physical-layer-processing}

Ye et al.~\cite{YeLiJuang2018OFDMDL} demonstrated that deep neural networks can jointly perform channel estimation and signal detection in OFDM systems, achieving near-optimal performance with lower complexity than separate estimation and detection stages. Their key insight was that the DNN implicitly learns the channel's statistical structure, eliminating the need for explicit pilot-based estimation.

Samuel et al.~\cite{SamuelDiskinWunder2019MIMODetect} proposed DetNet, an unfolded projected gradient descent network for MIMO detection. By unrolling iterative optimization into a fixed number of neural network layers, DetNet achieves near-maximum-likelihood detection performance with \(O(N_t^2)\) complexity instead of the exponential complexity of exhaustive ML search. This demonstrates the power of incorporating domain knowledge (the iterative structure of the detection problem) into the network architecture.

Dorner et al.~\cite{Dorner2018OverTheAirDL} proposed treating the entire communication system: modulator, channel, and demodulator: as an autoencoder, trained end-to-end to minimize bit error rate. This approach jointly optimizes all components, potentially discovering novel modulation and coding schemes that outperform separately designed modules. The autoencoder paradigm is conceptually related to the relay processing task: both involve learning an encoder-decoder pair that maps through a noisy channel.

Sun et al.~\cite{Sun2018LearnToOptimize} applied deep learning to the NP-hard problem of resource allocation in interference networks, demonstrating that a DNN can learn near-optimal power control policies orders of magnitude faster than conventional iterative algorithms.

In the context of relay-assisted communication systems, most existing works employ neural networks primarily for \emph{relay-selection} and \emph{resource-allocation}, particularly in decode-and-forward (DF) relaying scenarios, where learning-based classifiers or predictors are used to identify the optimal relay based on channel state information, signal-to-noise ratios, or energy constraints, while the underlying relay operation remains conventional \cite{Akdemir2024DLRelaySelection}.

A broader perspective on the role of machine learning in wireless communications is provided in \cite{Gunduz2019MLAir}, which highlights the potential of learning-based methods across multiple protocol layers, including cooperative communications.

Another related line of work exploits the structural similarities between amplify-and-forward (AF) relay networks and neural networks, treating relay nodes as neurons with nonlinear transfer functions and using deep-learning optimization tools to tune relay gains and exploit hardware-induced nonlinearities \cite{Bergel2023RelaysAsNeurons,BergelICASSP2024}. While these approaches demonstrate significant gains over classical linear AF schemes, they focus on analog gain optimization rather than symbol-level detection.

Classical cooperative communication frameworks based on AF and DF relaying remain well understood and widely adopted \cite{Nosratinia2004CooperativeCommunication}, but they rely on fixed linear forwarding or hard-decision rules. Consequently, despite substantial progress, relatively little attention has been devoted to learning the relay's symbol-level input-output mapping itself, motivating neural network-based detection at the relay as a generalized, data-driven alternative to conventional AF and DF processing.

\subsection{Neural Network Relay Processing}\label{sec:neural-network-relay-processing}

For relay processing specifically, a supervised learning approach trains a neural network \(f_\theta\) to minimize the empirical mean-squared error: (Equation~\ref{eq:mse-loss})

\begin{equation}\protect\phantomsection\label{eq:mse-loss}{\mathcal{L}(\theta) = \frac{1}{N} \sum_{i=1}^{N} (\hat{x}_i - x_i)^2, \quad \hat{x}_i = f_\theta(y_{i-w:i+w})
}\end{equation}
\addcontentsline{equ}{listequations}{\protect\numberline{\theequation}\quad mse-loss}

where \(\hat{x}_i = f_\theta(y_{i-w:i+w})\) is the network output based on a sliding window of \(2w+1\) received symbols, and \(x_i\) is the clean transmitted symbol. This window-based approach provides temporal context that enables the network to exploit statistical dependencies in the noise-corrupted signal.

A well-known example of \(f_\theta(\cdot)\) is the \emph{multilayer perceptron (MLP)}, a feedforward artificial neural network composed of an input layer, one or more fully connected hidden layers with nonlinear activation functions, and an output layer, trained using backpropagation to model nonlinear input--output relationships.
The MLP equation is shown in equation \ref{eq:mlp-relay}:

\begin{equation}\protect\phantomsection\label{eq:mlp-relay}{\mathbf{h} = \text{ReLU}(\mathbf{W}_1 \mathbf{w} + \mathbf{b}_1), \quad \hat{x} = \tanh(\mathbf{W}_2 \mathbf{h} + \mathbf{b}_2)
}\end{equation}
\addcontentsline{equ}{listequations}{\protect\numberline{\theequation}\quad mlp-relay}

The sliding window formulation ($i \in [i-w:i+w]$) is motivated by the observation that adjacent received symbols share common noise statistics and channel conditions. 
%While AWGN noise is i.i.d.\ across symbols (making the window unnecessary for a single-symbol estimator), the window provides the network with a local context that helps it learn a more robust decision boundary: effectively performing a form of implicit averaging that reduces variance. For fading channels, where adjacent symbols may experience correlated fading, the window becomes even more valuable. \AK{Didn't you assume white noise and i.i.d.\ Rayleigh fading? If so, why do you talk about correlation and how can $w>0$ help?}
Under the canonical model (i.i.d.\ fast fading, white noise), adjacent symbols are statistically independent, so the single received sample $y_R[i]$ is already a sufficient statistic for $x[i]$; a single-symbol estimator ($w=0$) therefore loses nothing and a window confers no theoretical benefit. The window is therefore not motivated by temporal correlation in the canonical setting; it is retained because it is required by the unknown-channel study of Chapter~\ref{sec:unknown-channels}, where a 3-tap ISI channel couples adjacent symbols and the relay must span that memory. Keeping a fixed windowed architecture across all experiments isolates the effect of the channel assumptions from the network structure. 

\textbf{Multi-SNR training} is a key design choice: by training on data generated at multiple SNR levels (5, 10, 15 dB in this thesis), the network learns a denoising function that generalizes across operating conditions rather than specializing to a single noise level. This trains a single estimator to work across a range of noise variances. In decision-theoretic terms the training objective averages the loss over the sampled SNRs, so what it minimises is \emph{Bayes} risk under the empirical SNR distribution, not the worst-case (minimax) risk over that range: nothing in the objective protects the SNR at which the network happens to do worst.

A critical question in applying neural networks to relay processing is the relationship between model complexity and performance. Prior work has generally assumed that larger models yield better performance, but this assumption has not been rigorously tested in the relay communication context. The bias--variance framework predicts that for a low-complexity target function like relay denoising, performance should plateau or degrade beyond a modest model size: a prediction investigated empirically in this thesis. Understanding this relationship is essential for practical deployment, where relay nodes may have limited computational resources.

\section{Generative Models for Signal Processing}\label{sec:generative-models-for-signal-processing}

Generative models offer an alternative paradigm for relay signal processing. Rather than directly learning a denoising function (discriminative approach), generative models learn the distribution of clean signals \(p(\mathbf{x})\) and use this knowledge to reconstruct the transmitted signal from noisy observations via Bayes' rule:
\begin{equation}\protect\phantomsection\label{eq:bayes-rule}{p(\mathbf{x} | \mathbf{y}) \propto p(\mathbf{y} | \mathbf{x}) p(\mathbf{x})}
\end{equation}

The generative approach is theoretically appealing because it separates the modeling of the signal prior from the noise model, potentially enabling better generalization to unseen noise conditions.

\subsection{Variational Autoencoders}\label{sec:variational-autoencoders}

\textbf{Variational Autoencoders (VAEs)} \cite{KingmaWelling2014VAE} learn a latent representation by maximizing the evidence lower bound (ELBO). The generative model posits that data \(\mathbf{x}\) is generated from a latent variable \(\mathbf{z} \sim p(\mathbf{z}) = \mathcal{N}(\mathbf{0}, \mathbf{I})\) through a decoder \(p_\theta(\mathbf{x} | \mathbf{z})\). Since the true posterior \(p(\mathbf{z} | \mathbf{x})\) is intractable, an encoder network \(q_\phi(\mathbf{z} | \mathbf{x})\) approximates it. The ELBO objective is derived from the log-evidence decomposition: (Equation~\ref{eq:vae-elbo})

\begin{equation}\protect\phantomsection\label{eq:vae-elbo}{\log p_\theta(\mathbf{x}) = \underbrace{\mathbb{E}_{q_\phi}\left[\log \frac{p_\theta(\mathbf{x}, \mathbf{z})}{q_\phi(\mathbf{z} | \mathbf{x})}\right]}_{\text{ELBO}(\theta, \phi; \mathbf{x})} + \underbrace{D_{KL}(q_\phi(\mathbf{z} | \mathbf{x}) \| p_\theta(\mathbf{z} | \mathbf{x}))}_{\geq 0}
}\end{equation}
\addcontentsline{equ}{listequations}{\protect\numberline{\theequation}\quad vae-elbo}

Since the KL divergence is non-negative, the ELBO is a lower bound on the log-evidence. Maximizing the ELBO simultaneously trains the encoder to approximate the true posterior and the decoder to reconstruct the data. The ELBO decomposes into a reconstruction term and a regularization term:

\begin{equation}\protect\phantomsection\label{eq:vae-loss}{\mathcal{L}_{\text{VAE}} = \underbrace{\mathbb{E}_{q_\phi(\mathbf{z}|\mathbf{x})}[\log p_\theta(\mathbf{x}|\mathbf{z})]}_{\text{Reconstruction quality}} - \underbrace{\beta \cdot D_{KL}(q_\phi(\mathbf{z}|\mathbf{x}) \| p(\mathbf{z}))}_{\text{Latent space regularity}}
}\end{equation}
\addcontentsline{equ}{listequations}{\protect\numberline{\theequation}\quad vae-loss}

The reconstruction term encourages the decoder to accurately reproduce the input from the latent code, while the KL term regularizes the latent space to be close to the standard Gaussian prior \(p(\mathbf{z})\). The \(\beta\) parameter (\(\beta\)-VAE) \cite{Higgins2017betaVAE} in Equation~\ref{eq:vae-loss} controls the trade-off between reconstruction quality and latent space smoothness: \(\beta < 1\) prioritizes reconstruction (beneficial for the relay task where accurate signal recovery is paramount), while \(\beta > 1\) encourages disentangled latent representations.

The \emph{reparameterization-trick} enables gradient-based optimization through the stochastic sampling layer:
\begin{equation}\protect\phantomsection\label{eq:repremtrized-trick}
\mathbf{z} \sim q_\phi(\mathbf{z} | \mathbf{x}) = \mathcal{N}(\boldsymbol{\mu}, \text{diag}(\boldsymbol{\sigma}^2))
\end{equation}
Instead of sampling from equation \ref{eq:repremtrized-trick}, the sample is expressed as:
\begin{equation}\protect\phantomsection\label{eq:repremtrized-trick-resample}
\mathbf{z} = \boldsymbol{\mu} + \boldsymbol{\sigma} \odot \boldsymbol{\epsilon}
\end{equation}
\addcontentsline{equ}{listequations}{\protect\numberline{\theequation}\quad repremtrized-trick}

where \(\boldsymbol{\epsilon} \sim \mathcal{N}(\mathbf{0}, \mathbf{I})\).
This moves the stochasticity outside the computational graph, allowing backpropagation through the encoder.

For relay signal processing, the VAE learns a compressed latent representation of the clean signal manifold. During inference, the noisy received signal is encoded into the latent space, and the decoder maps it back to a denoised estimate. The regularized latent space provides implicit denoising: noisy inputs that map to regions far from the learned signal manifold are pulled back toward high-probability regions during decoding. Sampling \(\mathbf{z}\) stochastically at inference does, however, inject variance into an otherwise deterministic estimate, which is a plausible handicap for a deterministic BPSK denoising target; Section~\ref{sec:generative-vs.-discriminative-paradigms-for-relay-processing} reports what this is and is not found to cost in the measurements.

\subsection{Conditional Generative Adversarial Networks}\label{sec:conditional-generative-adversarial-networks}

\textbf{Conditional GANs (cGANs)} \cite{MirzaOsindero2014cGAN} learn a denoising map by adversarial training: a generator conditioned on the noisy observation produces a clean estimate, and a discriminator is trained to separate generated from true clean symbols, the two optimizing a minimax objective. The original formulation is unstable to train, because the divergence the discriminator implicitly minimizes saturates when the two distributions do not overlap and stops supplying useful gradients. The Wasserstein GAN with gradient penalty (WGAN-GP) \cite{Gulrajani2017WGANGP} replaces that objective with the Wasserstein distance, estimated through a critic constrained to be approximately 1-Lipschitz by penalizing the norm of its gradient at points interpolated between real and generated samples. This is the variant implemented in this thesis (Section~\ref{sec:relay-strategies}).

The cGAN is described here and in Chapter~\ref{sec:methods} for completeness, but it is \emph{excluded from every reported comparison} on cost grounds: training it took roughly two GPU-hours against the minimal MLP's few seconds on a CPU, which is not a defensible expenditure for a relay-denoising target this simple. No claim in this thesis rests on it, and the theory above is therefore kept deliberately brief.

\subsection{Generative vs.~Discriminative Paradigms for Relay Processing}\label{sec:generative-vs.-discriminative-paradigms-for-relay-processing}

The application of generative models to relay processing has, to the best of our knowledge, not been extensively studied. The key question is whether the generative inductive bias (learning the data distribution rather than just the input-output mapping) provides any advantage for the relay denoising task. For BPSK, the clean signal distribution is trivially simple (a discrete distribution on \(\{-1, +1\}\)), suggesting that the generative overhead may not be justified. This thesis compares VAE-based relay processing against classical and supervised learning methods on a common benchmark, and the results show that the generative paradigm provides no significant advantage for this particular task. A conditional GAN relay is implemented and described in Chapter~\ref{sec:methods} but is excluded from the reported comparison on cost grounds, so no claim is made about it here. On the canonical setup the VAE tracks the supervised MLP and symbol-wise DF rather than trailing them, reaching $0.00972$ at 20~dB against DF's $0.00972$ and the MLP's $0.00992$ (Table~\ref{tbl:table2}), and it stays inside the feedforward group at every SNR in the equal-parameter comparison as well (Table~\ref{tbl:table8}). The generative bias is therefore neither an advantage nor a handicap here: it costs parameters and training time to arrive where a 169-parameter regressor already is. One qualification remains, unchanged by this: the relay samples $\mathbf{z}\sim q_\phi(\mathbf{z}\mid\mathbf{x})$ stochastically at inference, and the deterministic variant that substitutes the posterior mean $\boldsymbol{\mu}$ at test time was not evaluated, so nothing here bears on whether that variant would do better.

\section{Sequence Models: Transformers, State Space Models}\label{sec:sequence-models-transformers-and-state-space-models}

Recent advances in sequence modeling have produced two competing paradigms with distinct computational properties and inductive biases. Both have achieved state-of-the-art results in natural language processing, but their relative merits for physical-layer signal processing (where sequences are real-valued temporal signals rather than discrete tokens) have, to the best of our knowledge, not been previously investigated.

\subsection{Transformers and the Attention Mechanism}\label{sec:transformers-and-the-attention-mechanism}

\textbf{Transformers} \cite{Vaswani2017Attention} use multi-head self-attention to capture global dependencies in sequences. The core attention mechanism computes a weighted combination of value vectors, where the weights are derived from the compatibility of query and key vectors:

\begin{equation}\protect\phantomsection\label{eq:attention}{\text{Attention}(\mathbf{Q}, \mathbf{K}, \mathbf{V}) = \text{softmax}\left(\frac{\mathbf{Q}\mathbf{K}^T}{\sqrt{d_k}}\right)\mathbf{V}
}\end{equation}
\addcontentsline{equ}{listequations}{\protect\numberline{\theequation}\quad attention}

where \(\mathbf{Q}, \mathbf{K} \in \mathbb{R}^{n \times d_k}\) and \(\mathbf{V} \in \mathbb{R}^{n \times d_v}\) are obtained from the input via learned linear projections \(W^Q, W^K, W^V\). The scaling factor \(1/\sqrt{d_k}\) prevents the dot products from growing too large in magnitude, which would push the softmax into saturation regions with vanishing gradients.

\textbf{Multi-head attention} extends this by running \(h\) parallel attention heads, each with independent projections, and concatenating their outputs: (Equation~\ref{eq:multihead-attention})
\begin{equation}\protect\phantomsection\label{eq:multihead-attention}{\text{MultiHead}(\mathbf{X}) = \text{Concat}(\text{head}_1, \dots, \text{head}_h) W^O,\\
\text{head}_i = \text{Attention}(\mathbf{X} W_i^Q, \mathbf{X} W_i^K, \mathbf{X} W_i^V)
}\end{equation}
\addcontentsline{equ}{listequations}{\protect\numberline{\theequation}\quad multihead-attention}

Each head can attend to different aspects of the input: for signal processing, one head might focus on the immediate neighbors (local denoising) while another captures longer-range patterns. The total parameter count for multi-head attention is:

\begin{equation}\protect\phantomsection\label{eq:multihead-attention-parameters-count}{3 d_{\text{model}}^2 + d_{\text{model}}^2 = 4 d_{\text{model}}^2}
\end{equation}
\addcontentsline{equ}{listequations}{\protect\numberline{\theequation}\quad multihead-attention-parameters-count}

for \(Q, K, V\) projections and the output projection. The attention matrix:
\begin{equation}\protect\phantomsection\label{eq:multihead-attention-matrix}{\mathbf{A} = {softmax({\frac{\mathbf{Q}\mathbf{K}^T}{\sqrt{d_k}}})} \in \mathbb{R}^{n \times n}}
\end{equation}
\addcontentsline{equ}{listequations}{\protect\numberline{\theequation}\quad multihead-attention-matrix}

computes pairwise interactions between all \(n\) positions, yielding a time and memory complexity of: $O(n^2)$. For the 11-symbol relay window used in this thesis, this is negligible (\(11^2 = 121\) entries). However, the quadratic cost becomes prohibitive for longer sequences (e.g., \(n = 1000\) symbols), motivating linear-time alternatives.

\textbf{Positional encoding} is necessary because the attention mechanism is permutation-equivariant (it treats input positions symmetrically). This thesis uses sinusoidal positional encoding: (Equation~\ref{eq:positional-encoding})

\begin{equation}\protect\phantomsection\label{eq:positional-encoding}{PE_{(pos, 2k)} = \sin(\frac{pos}{10000^{2k/d}}), \quad PE_{(pos, 2k+1)} = \cos(\frac{pos}{10000^{2k/d}})
}\end{equation}
\addcontentsline{equ}{listequations}{\protect\numberline{\theequation}\quad positional-encoding}

which injects position information into the embeddings, enabling the model to distinguish between symbols at different positions in the window.

\subsection{Structured State Space Models}\label{sec:structured-state-space-models}

Structured state space models (SSMs) \cite{GuGoelRe2022SSM} treat a sequence model as a discretized linear dynamical system. A continuous system $\dot{\mathbf{x}}(t)=\mathbf{A}\mathbf{x}(t)+\mathbf{B}u(t)$, $y(t)=\mathbf{C}\mathbf{x}(t)+Du(t)$ is discretized with a step size $\Delta$ (zero-order hold, $\bar{\mathbf{A}}=\exp(\Delta\mathbf{A})$) to give the recurrence

\begin{equation}\protect\phantomsection\label{eq:ssm-discrete}{\mathbf{x}_k = \bar{\mathbf{A}} \mathbf{x}_{k-1} + \bar{\mathbf{B}} u_k, \quad y_k = \mathbf{C} \mathbf{x}_k + D u_k
}\end{equation}
\addcontentsline{equ}{listequations}{\protect\numberline{\theequation}\quad ssm-discrete}

which is a linear recursive filter with a learned state transition. Two ingredients make this competitive with attention. The HiPPO framework supplies a principled initialization for $\mathbf{A}$ under which the state maintains a compressed polynomial approximation of the input history, and S4 \cite{GuGoelRe2022SSM} shows that a structured (diagonal or diagonal-plus-low-rank) $\mathbf{A}$ makes the recurrence cheap to evaluate. The resulting model is linear-time in sequence length, in contrast to attention's quadratic cost.

\subsection{Mamba and Mamba-2}\label{sec:mamba-selective-state-spaces}

\textbf{Mamba} \cite{gu2023mamba} makes the SSM parameters \emph{input-dependent}, turning the time-invariant system into a time-varying one:

\begin{equation}\protect\phantomsection\label{eq:mamba-selective}{\Delta_k = \text{softplus}(W_\Delta u_k + b_\Delta), \quad \mathbf{B}_k = W_B u_k, \quad \mathbf{C}_k = W_C u_k
}\end{equation}
\addcontentsline{equ}{listequations}{\protect\numberline{\theequation}\quad mamba-selective}

With $\mathbf{A}$ diagonal and negative, $\Delta_k$ acts as a learned gate: a large $\Delta_k$ drives $\bar{\mathbf{A}}_k=\exp(\Delta_k\mathbf{A})$ toward zero, resetting the state so the new input dominates, while a small $\Delta_k$ leaves $\bar{\mathbf{A}}_k\approx\mathbf{I}$, preserving the state and ignoring the input. This selectivity is the property of interest for relay processing: in principle the model can retain signal structure while suppressing noise, deciding per sample which of the two an observation carries. The selective layer is wrapped in a gated block with an expand-factor-2 projection and a short causal convolution.

\begin{remark}
These adaptive-filtering connections become directly relevant only in the unknown-channel study of Chapter~\ref{sec:unknown-channels}; on the canonical memoryless channel the relay reduces to a per-symbol function and there is no temporal structure for a state to track.
\end{remark}

\begin{figure}[htbp]
    \centering
    \includegraphics[width=\linewidth]{results/mamba_ssd_s6.png}
    \caption[Mamba: selective state spaces]{Structured SSMs map each channel independently through a higher-dimensional latent state; the selective variant makes the transition input-dependent so the state is reset or held per sample}
    \label{fig:mamba_ssd_s6}
\end{figure}

\textbf{Mamba-2 (SSD)} \cite{DaoGu2024TransformersSSM} reformulates the same computation through an algebraic duality between the linear recurrence and a structured matrix multiplication: the output can be written $\mathbf{y}=M\cdot(\mathbf{B}\odot\mathbf{u})$ for a lower-triangular semi-separable matrix $M$. The practical consequence is that the sequence can be processed in chunks, with a batched matrix multiply inside each chunk and a single state hand-off between chunks, so the number of sequential steps falls from one per sample to one per chunk. Mamba-2 also forms the state-space parameters at the start of the block rather than sequentially along it. Both variants are evaluated here; at the short relay windows this thesis uses, the distinction is architectural rather than a source of measured advantage.

\subsection{State Space Models for Signal Processing}\label{sec:state-space-models-for-signal-processing}

The application of state space models to physical-layer signal processing is, to the best of our knowledge, largely unexplored. Classical signal processing relies extensively on LTI state space models (Kalman filters, Wiener filters), and the SSM framework provides a natural neural extension of these classical tools. The connection is direct: a trained SSM with fixed (input-independent) parameters implements a learnable linear filter, while the selective Mamba variant implements a learnable \emph{adaptive} filter. For the relay denoising task, this means Mamba can potentially learn the optimal filter structure from data, without requiring manual specification of filter order, cutoff frequencies, or adaptation algorithms.

This thesis presents, to the best of our knowledge, the first comparison of Mamba-S6, Mamba-2 SSD, and Transformers for relay communication. The relay window used throughout is short ($n=11$), which is precisely the regime where attention's quadratic cost is negligible and a linear-time recurrence has nothing to recover; the throughput advantages reported for the SSD formulation at longer context lengths \cite{DaoGu2024TransformersSSM} are therefore not exercised by this application and are not measured here.

\section{Theoretical Foundations: Channel Models}\label{sec:theoretical-foundations-channel-models}

This section presents the theoretical BER analysis for each channel model used in this study, derives the closed-form expressions, and validates them against Monte Carlo simulation. The theoretical-simulative comparison serves two purposes: (i) it verifies the correctness of the simulation framework, and (ii) it establishes the baseline performance that AI relays must beat.

\subsection{AWGN Channel: Theoretical Analysis}\label{sec:awgn-channel-theoretical-analysis}

The AWGN channel adds zero-mean Gaussian noise to the transmitted signal: (Equation~\ref{eq:awgn-channel})

\begin{equation}\protect\phantomsection\label{eq:awgn-channel}{y = x + n, \quad n \sim \mathcal{N}(0, \sigma^2), \quad \sigma^2 = \frac{P_s}{\text{SNR}_{\text{linear}}}
}\end{equation}
\addcontentsline{equ}{listequations}{\protect\numberline{\theequation}\quad awgn-channel}

where \(P_s = \mathbb{E}[|x|^2] = 1\) for unit-power BPSK symbols. The noise variance is inversely proportional to the linear SNR.

\subsubsection{Single-hop BER}\label{sec:awgn-single-hop-ber}
 For BPSK \(\pm 1\) symbols and real-valued AWGN with variance \(\sigma^2 = 1/\text{SNR}\), an error occurs when the noise pushes the received sample past the decision boundary at the origin. The theoretical BER is: (Equation~\ref{eq:awgn-ber})

\begin{equation}\protect\phantomsection\label{eq:awgn-ber}{P_e^{\text{AWGN}} = Q\!\left(\frac{1}{\sigma}\right) = Q\left(\sqrt{\text{SNR}}\right) = \frac{1}{2}\,\text{erfc}\!\left(\sqrt{\frac{\text{SNR}}{2}}\right)
}\end{equation}
\addcontentsline{equ}{listequations}{\protect\numberline{\theequation}\quad awgn-ber}

where \(Q(x) = \frac{1}{2}\text{erfc}(x/\sqrt{2})\) is the Gaussian Q-function and \(\text{SNR} = P_s / \sigma^2\) is the ratio of signal power to noise variance.

\begin{remark}[Noise convention]\label{noise-convension}
The standard textbook result \(P_e = Q(\sqrt{2 E_b/N_0})\) assumes that the one-sided noise PSD is \(N_0\), giving a baseband noise variance of \(N_0/2\) per real dimension. In our AWGN implementation the noise is purely real with variance \(\sigma^2 = P_s/\text{SNR}\), so \(N_0 = 2\sigma^2 = 2P_s/\text{SNR}\) and \(E_b/N_0 = \text{SNR}/2\). Substituting yields \(Q(\sqrt{2 \cdot \text{SNR}/2}) = Q(\sqrt{\text{SNR}})\), which is the expression used throughout this work.

By contrast, the fading channels add \textbf{complex} Gaussian noise with total power \(1/\text{SNR}\) (i.e.~\(\sigma_I^2 = \sigma_Q^2 = 1/(2\,\text{SNR})\)) and extract the real part after ZF equalization, halving the effective noise variance per decision dimension. This naturally introduces a factor of 2 inside the Q-function argument for all fading-channel BER expressions (the fading-channel analysis above), making them consistent with the standard textbook forms.
\end{remark}

\subsubsection{Two-hop AF BER}\label{sec:awgn-two-hop-af-ber}
For amplify-and-forward with equal-SNR hops, the effective end-to-end SNR is:

\begin{equation}\protect\phantomsection\label{eq:af-snr-eff}{\text{SNR}_{\text{eff}}^{\text{AF}} = \frac{\text{SNR}_1 \cdot \text{SNR}_2}{\text{SNR}_1 + \text{SNR}_2 + 1} = \frac{\gamma^2}{2\gamma + 1}
}\end{equation}
\addcontentsline{equ}{listequations}{\protect\numberline{\theequation}\quad af-snr-eff}

where
\(\gamma = \text{SNR}\) is the per-hop SNR. The resulting BER is

\begin{equation}\protect\phantomsection\label{eq:af-p-eff}{P_e^{\text{AF}} = Q\left(\sqrt{\text{SNR}_{\text{eff}}}\right)
}\end{equation}
\addcontentsline{equ}{listequations}{\protect\numberline{\theequation}\quad af-p-eff}

under the same real-noise convention as the single-hop case (Remark~\ref{noise-convension}). At high SNR,

\begin{equation}\protect\phantomsection\label{eq:gamma-high-snr}{\text{SNR}_{\text{eff}} \approx \frac{\gamma}{2}}
\end{equation}
\addcontentsline{equ}{listequations}{\protect\numberline{\theequation}\quad gamma-high-snr}

confirming the well-known 3 dB penalty of AF relaying.

\subsubsection{Two-hop DF BER}\label{sec:awgn-two-hop-df-ber}
For a decode-and-forward relay with equal-SNR hops, an error occurs at the destination if exactly one hop introduces an error:

\begin{equation}\protect\phantomsection\label{eq:df-ber-awgn}{P_e^{\text{DF}} = P_{e,1} + P_{e,2} - 2P_{e,1}P_{e,2}
}\end{equation}
\addcontentsline{equ}{listequations}{\protect\numberline{\theequation}\quad df-ber-awgn}

With equal hops (i.e.\ \(P_{e,1} = P_{e,2} = P_e^{\text{AWGN}}\)), this becomes

\begin{equation}\protect\phantomsection\label{eq:equal-hop-snr}{P_e^{\text{DF}} = 2P_e(1 - P_e)}
\end{equation}\addcontentsline{equ}{listequations}{\protect\numberline{\theequation}\quad equal-hop-snr}

At high SNR, \(P_e \ll 1\) and
\begin{equation}\protect\phantomsection\label{eq:df-high-snr-two-hop-penalty}{P_e^{\text{DF}} \approx 2P_e}
\end{equation}\addcontentsline{equ}{listequations}{\protect\numberline{\theequation}\quad df-high-snr-two-hop-penalty}
i.e., the two-hop penalty is approximately a factor of 2 in BER.

\subsection{Rayleigh Fading Channel: Theoretical Analysis}\label{sec:rayleigh-fading-channel-theoretical-analysis}

The Rayleigh fading channel models non-line-of-sight (NLOS) propagation where the signal undergoes multiplicative fading:

\begin{equation}\protect\phantomsection\label{eq:rayleigh-channel}{y = hx + n, \quad h \sim \mathcal{CN}(0, 1), \quad n \sim \mathcal{CN}(0, \sigma^2)
}\end{equation}
\addcontentsline{equ}{listequations}{\protect\numberline{\theequation}\quad rayleigh-channel}

The fading coefficient \(h\) is a circularly-symmetric complex Gaussian random variable, so its magnitude \(|h|\) follows a Rayleigh distribution:

\begin{equation}\protect\phantomsection\label{eq:rayleigh-pdf}{f_{|h|}(r) = 2r \cdot e^{-r^2}, \quad r \geq 0
}\end{equation}
\addcontentsline{equ}{listequations}{\protect\numberline{\theequation}\quad rayleigh-pdf}

with \(\mathbb{E}[|h|^2] = 1\) (unit average power). The instantaneous SNR after equalization (\(\hat{x} = y/h\)) becomes \(\gamma_h = |h|^2 \cdot \text{SNR}\), which is exponentially distributed.

\subsubsection{Single-hop BER}\label{sec:rayleigh-single-hop-ber}
Averaging the conditional BER \(P_e(\gamma_h) = Q(\sqrt{2\gamma_h})\) over the exponential distribution of \(\gamma_h\) yields the closed-form \cite[Eq.~14-4-15]{ProakisSalehi2008DigitalComm} (Equation~\ref{eq:rayleigh-ber}):

\begin{equation}\protect\phantomsection\label{eq:rayleigh-ber}{P_e^{\text{Rayleigh}} = \frac{1}{2}\left(1 - \sqrt{\frac{\bar{\gamma}}{1 + \bar{\gamma}}}\right)
}\end{equation}
\addcontentsline{equ}{listequations}{\protect\numberline{\theequation}\quad rayleigh-ber}

where \(\bar{\gamma} = \text{SNR}\) is the average SNR. At high SNR (\(\bar{\gamma} \gg 1\)):

\begin{equation}\protect\phantomsection\label{eq:rayleigh-ber-approx}{P_e^{\text{Rayleigh}} \approx \frac{1}{4\bar{\gamma}}
}\end{equation}
\addcontentsline{equ}{listequations}{\protect\numberline{\theequation}\quad rayleigh-ber-approx}

This \(1/\text{SNR}\) decay is fundamentally slower than the exponential decay of AWGN (\(Q(\sqrt{\gamma}) \sim e^{-\gamma/2}\)), explaining why Rayleigh fading is significantly more challenging. The channel is \textbf{diversity-limited}: deep fades (where \(|h| \approx 0\)) cause errors regardless of the average SNR. Combating this diversity limitation is the primary motivation for the relay architectures studied in this thesis.

\subsubsection{Two-hop DF BER}\label{sec:rayleigh-two-hop-df-ber}
The two-hop DF relay over Rayleigh fading follows the same composition rule as AWGN:

\begin{equation}\protect\phantomsection\label{eq:df-ber-rayleigh}{P_e^{\text{DF,Rayleigh}} = 2P_e^{\text{Rayleigh}}(1 - P_e^{\text{Rayleigh}})
}\end{equation}
\addcontentsline{equ}{listequations}{\protect\numberline{\theequation}\quad df-ber-rayleigh}

\section{Channel Coding: Convolutional Codes and Trellis Decoding}\label{sec:channel-coding-convolutional-codes}

The channel models above assume an uncoded link. Chapter~\ref{sec:experiments}'s supplementary coded study (Section~\ref{sec:coded-block-df}) departs from that assumption, adding a convolutional code and a block-DF relay that decodes it. This section gives the minimal background needed to follow that study: what a convolutional code's rate and constraint length mean, how a receiver recovers the transmitted sequence from a noisy one, and how a single code is stretched across several rates.

\subsection{Convolutional Encoding: Rate and Constraint Length}\label{sec:convolutional-encoding-rate-constraint-length}

A convolutional encoder maps each group of $k$ input bits to $n$ coded output bits using a shift register of \emph{constraint length} $K$ -- the number of $k$-bit input groups (including the current one) that influence each output. The code \emph{rate} \cite{ProakisSalehi2008DigitalComm} is
\begin{equation}\protect\phantomsection\label{eq:conv-code-rate}{R = \frac{k}{n}}\end{equation}
\addcontentsline{equ}{listequations}{\protect\numberline{\theequation}\quad conv-code-rate}

so a rate-$1/2$ code, the mother code used throughout Chapter~\ref{sec:experiments}, emits two coded bits per information bit. This thesis uses $k=1$ throughout, one information bit per trellis step; the encoder's output then depends on the current bit and the $K-1$ preceding ones held in the shift register, so its evolution is a finite-state machine with \cite{ProakisSalehi2008DigitalComm}
\begin{equation}\protect\phantomsection\label{eq:conv-trellis-states}{S = 2^{K-1}}\end{equation}
\addcontentsline{equ}{listequations}{\protect\numberline{\theequation}\quad conv-trellis-states}

states, and the resulting state diagram unrolled over time is the code's \emph{trellis}: every valid coded sequence corresponds to exactly one path through it. A larger $K$ gives the code more memory to spread redundancy across, at the cost of a larger trellis to search; Chapter~\ref{sec:experiments} measures this trade-off directly by sweeping $K \in \{3, 5, 7\}$, i.e.\ $4$, $16$ and $64$ states (Section~\ref{sec:coded-block-df}). A frame is conventionally \emph{zero-tail terminated}: $K-1$ zero bits are appended after the information bits to drive the encoder back to the all-zero state, so the trellis both starts and ends at a known point rather than an unknown one.

\subsection{Maximum-Likelihood Sequence Decoding: The Viterbi Algorithm}\label{sec:viterbi-ml-decoding}

Because every valid coded sequence is one path through the trellis, recovering the most likely transmitted sequence from a noisy received one is a shortest-path problem: find the trellis path whose hypothesized coded symbols are collectively closest to what was received. A brute-force search is infeasible -- an $L$-bit information frame admits $2^{L}$ trellis paths. The Viterbi algorithm nonetheless solves the problem exactly, in time linear in $L$, by a dynamic-programming recursion over the trellis states: at each step it extends every surviving path by one branch and keeps only the lowest-cost path into each of the $S$ states (add-compare-select), so the cost per step depends on $S$ rather than on the number of paths \cite{Forney1972MLSE, ProakisSalehi2008DigitalComm}. For a memoryless channel with an appropriate branch metric -- squared Euclidean distance to the received sample under AWGN, as used by Chapter~\ref{sec:experiments}'s coded block-DF relay \cite{Forney1972MLSE} -- the recovered sequence is
\begin{equation}\protect\phantomsection\label{eq:viterbi-ml-metric}{\hat{c} = \operatorname*{arg\,min}_{c \in \mathcal{C}} \sum_{k} \left| y_k - x(c_k) \right|^2}\end{equation}
\addcontentsline{equ}{listequations}{\protect\numberline{\theequation}\quad viterbi-ml-metric}

where $\mathcal{C}$ is the set of trellis-valid coded sequences and $x(c_k)$ the modulated symbol for coded bits $c_k$ -- exactly the maximum-likelihood sequence estimate for the channel, with no approximation. Optimality is over the sequence, however, not over each bit independently, and this shapes how the decoder fails. When noise makes an incorrect path cheaper than the transmitted one, that path typically diverges from the correct path for several trellis steps before merging back, so a single decoding failure corrupts a \emph{run} of information bits rather than an isolated one. Residual errors therefore arrive in bursts, and below the code's operating threshold these events become frequent enough that the burst penalty outweighs the coding gain. This is the mechanism behind the low-SNR coding penalty Chapter~\ref{sec:experiments} measures (Section~\ref{sec:coded-block-df}), where a fixed-rate code performs worse than no coding at all -- a memoryless per-symbol slicer has no trellis to diverge along, and so cannot fail this way.

\subsection{Soft-Decision Decoding: The Forward-Backward (BCJR) Algorithm}\label{sec:soft-decision-bcjr-decoding}

Viterbi decoding returns one hard path through the trellis and discards how close the alternatives were. The BCJR algorithm \cite{BahlCockeJelinekRaviv1974BCJR} instead computes, for every trellis transition, its posterior probability given the \emph{entire} received sequence, using both a forward recursion $\alpha_k(s)$ (the likelihood of reaching state $s$ at step $k$ from everything received so far) and a backward recursion $\beta_k(s)$ (the likelihood of everything still to come, given state $s$ at step $k$):
\begin{equation}\protect\phantomsection\label{eq:bcjr-posterior}{P(s_{k-1}=s',\, s_k=s \mid \mathbf{y}) \propto \alpha_{k-1}(s')\,\gamma_k(s',s)\,\beta_k(s)}\end{equation}
\addcontentsline{equ}{listequations}{\protect\numberline{\theequation}\quad bcjr-posterior}

where $\gamma_k(s',s)$ is the branch's transition-and-observation probability. Where Viterbi minimizes the probability of choosing the wrong \emph{sequence}, BCJR minimizes the error probability of each \emph{symbol} individually \cite{BahlCockeJelinekRaviv1974BCJR}; the two optimize different criteria, which is why neither dominates the other in Chapter~\ref{sec:experiments}'s measurements. Marginalizing these joint probabilities over the transitions consistent with each coded bit or symbol gives a \emph{soft} estimate -- a posterior mean or log-likelihood ratio -- in place of Viterbi's hard decision. This is the read-out Chapter~\ref{sec:experiments} calls \emph{soft block-DF} (Section~\ref{sec:coded-block-df}): the relay forwards the BCJR posterior mean instead of committing to the single most likely codeword.

\subsection{Puncturing, Spectral Efficiency, and Rate Adaptation}\label{sec:puncturing-spectral-efficiency}

A code's rate and the modulation's constellation size $M$ together set how much information it carries per channel use, its spectral efficiency \cite{ProakisSalehi2008DigitalComm}:
\begin{equation}\protect\phantomsection\label{eq:spectral-efficiency}{\eta = R \cdot \log_2 M \quad \text{[bits/symbol]}}\end{equation}
\addcontentsline{equ}{listequations}{\protect\numberline{\theequation}\quad spectral-efficiency}

so a rate-$1/2$ code halves the information rate of whatever modulation carries it, and comparing a coded link against an uncoded one at the same SNR without accounting for this is not a like-for-like comparison (Section~\ref{sec:coded-latency-throughput} revisits Chapter~\ref{sec:experiments}'s coded-versus-uncoded comparison on this basis). Rates other than a code's native $k/n$ are obtained by \emph{puncturing}: deleting a fixed, periodic subset of the encoder's output bits according to a puncturing pattern, and at the decoder re-inserting those positions as erasures (soft value zero) before running the identical trellis search -- so one Viterbi or BCJR decoder implementation serves every punctured rate from the same mother code. This is how Chapter~\ref{sec:experiments}'s rate-$1/2$ mother code is stretched to rates $2/3$ and $3/4$ for its link-adaptation study (Section~\ref{sec:coded-rate-adaptation}).

Spectral efficiency alone is not the quantity a link adapts on, because it counts bits that were transmitted rather than bits that arrived. A higher-rate choice buys more bits per symbol but raises the frame error rate, and any real protocol discards a frame with residual errors. The operative objective is therefore \emph{goodput}, the information rate actually delivered:
\begin{equation}\protect\phantomsection\label{eq:goodput}{G = \eta \cdot (1 - \mathrm{FER}) \quad \text{[information bits/symbol]}}\end{equation}
\addcontentsline{equ}{listequations}{\protect\numberline{\theequation}\quad goodput}

Maximizing $G$ over a grid of modulation-and-coding schemes at each SNR -- the standard adaptive modulation and coding (AMC) problem -- traces the link-adaptation envelope that Section~\ref{sec:coded-rate-adaptation} measures. The maximization is unconstrained: the FER penalty in Equation~\ref{eq:goodput} already prices in unreliability, so no separate reliability constraint is imposed.
